\documentclass[11pt]{article}
\usepackage[margin=0.8in]{geometry}
\usepackage{booktabs,hyperref}
\hypersetup{colorlinks=true,urlcolor=blue}
\emergencystretch=3em
\title{Impala performance experiment: Step 4}
\author{Batch independent chains into larger GPU calls}
\date{September 30, 2026}
\begin{document}\maketitle
\section*{Outcome}
Batching four chains completes faster than the historical serial GPU workflow,
but the current prototype does not beat independent compiled-CPU workers on
this machine. Use CPU parallelism for the tested pooled workflow. Keep GPU
batching experimental. The full four-chain posterior comparison passed
the predeclared quality gates. Production sampler code and defaults are unchanged;
backend cleanup, Mac testing and hierarchical algorithm changes were deferred.
\section*{How independence is preserved}
Each spawned process owns its sampler setup, adaptation and Python/NumPy random
state. Seeds remain 42, 123, 456 and 789. The existing pooled sampler runs
unchanged. When a chain needs predictions, it submits its material parameters
and waits. The coordinator concatenates one request from each active worker,
evaluates the combined batch, and splits the predictions back into the original
request order. It does not combine likelihoods, targets, random draws or
adaptation. Different request lengths and completion times are supported.

CPU-parallel workers evaluate their own model directly. A CPU-batched control
uses the same request coordinator, but evaluates the combined batch on CPU.
GPU-batched uses the existing float64 CUDA prototype. This separates batching
correctness, GPU acceleration and the practical multicore CPU alternative.
\section*{Accuracy validation}
The complete hardware-enabled suite passed all 94 tests in 115.67 seconds.
Three new boundary tests check heterogeneous request sizes, row ordering,
invalid requests, predictions and weighted likelihoods. Prediction tolerances
remain relative $10^{-8}$ and absolute $10^{-10}$ Mbar; likelihood absolute error
must be at most $10^{-6}$. The original numerical fixture was not overwritten.

In the four-chain, 3,000-state, 50-temperature comparisons, all temperature-chain
parameter arrays and recorded counters matched the independent CPU workers
exactly for both batched CPU and batched GPU. Noise differences passed relative
$10^{-8}$ and absolute $10^{-15}$ gates. Prediction differences passed the
original tolerances. The eight-chain comparison likewise matched cold-chain
parameters and counters exactly. Batching can introduce rounding even on CPU,
because interpolation uses flattened strain coordinates with experiment offsets.
\section*{End-to-end timing}
All figures below are elapsed seconds until every chain has finished sampling.
Warm-up, JIT and output compression are excluded and startup is recorded
separately. Transfers, IPC, synchronization and host numerical work are included.
A postprocessing barrier prevents completed workers from writing outputs or
calculating predictions while another chain is still sampling. Backends were
run sequentially. All chains use 50 temperatures and 100 history steps.
\begin{center}\begin{tabular}{rrrrr}\toprule
Chains & States/chain & Parallel CPU (s) & Batched GPU (s) & Speedup\\\midrule
4 & 3,000 & 18.53 & 23.99 & 0.77\\
8 & 3,000 & 20.59 & 39.26 & 0.52\\
4 & 40,000 & 252.64 & 334.53 & 0.76\\
\bottomrule\end{tabular}\end{center}
Speedup is CPU wall time divided by GPU wall time; values below one mean slower.
These are measured runs on this machine, not timing confidence intervals.
The eight-chain probe adds seeds 111, 222, 333 and 444; it is a scaling probe,
not a convergence study. The CPU has 24 cores (Intel Core Ultra 9 285HX), with
one numerical-library thread per worker. The GPU is the same RTX PRO 5000
Blackwell laptop GPU used in step 3.

For four short chains, the single CPU coordinator took 46.24 s, versus 23.99 s
with GPU evaluation. That comparison shows acceleration of the combined work,
but the independent CPU workers took only 18.53 s. Actual mean combined batches
were 159.1 parameter sets for four chains and 317.3 for eight, below the nominal
200 and 400 because proposals are filtered by constraints. Their corresponding
GPU evaluation totals were 18.11 and 30.21 s; total sampling wall times were
23.99 and 39.26 s. Host work, transfers and synchronization remain material costs.
\section*{Full posterior-quality comparison}
Four independent chains per backend: 40,000 states, discard the first 20,000,
no thinning, 50 temperatures. The same predeclared step-3 gates apply: all 16
parameter/noise quantities need rank R-hat below 1.01, bulk/tail ESS at least 400,
and MCSE(mean)/SD below 0.05. Posterior means must agree within four combined
MCSEs and 0.1 pooled posterior SD; prediction means within 0.1 experiment noise
SD. Predictions use the same 200 retained draws per chain as the baseline.
\begin{center}\begin{tabular}{lrr}\toprule
Diagnostic & Parallel CPU & Batched GPU\\\midrule
Maximum R-hat & 1.0026 & 1.0026\\
Minimum bulk ESS & 2253.5 & 2253.5\\
Minimum tail ESS & 3345.3 & 3345.3\\
Minimum bulk ESS/s & 8.9196 & 6.7363\\
\bottomrule\end{tabular}\end{center}
Overall gates passed. Maximum mean difference: 1.73e-12 combined MCSEs and
8.79e-15 posterior SD; maximum prediction-mean difference: 7.33e-13 noise SD.
Exact parameter trajectories for each GPU seed: 42: True, 123: True, 456: True, 789: True.
The independent CPU runs reproduce the original saved baseline arrays: True.
Detailed differences, including noise arrays and interval endpoints, are in JSON
and CSV artifacts. The hierarchical mixing problem is outside this experiment.

ESS/s uses the wall time for all four chains together, not the sum of the
individual worker times. The GPU/CPU ESS/s ratios range from 0.755 to 0.755.
This is a current multicore CPU comparison; it must not be substituted with
historical timings from serial execution. Quality and speed are separate gates:
accurate posterior draws do not establish a GPU performance benefit.
\section*{Tracking, limitations and reproduction}
Branch: \texttt{perf/gpu-batched-chains}. Initial code: \texttt{5588570}.\\
Long-run code: \texttt{3e1ad43}. Checkpoint: \texttt{perf-step4-20260930}.
The CPU completion timeout was corrected between runs, with the suite controller
paused, then smoke-tested. It affects driver liveness only. Initial short runs
record their source hashes; long runs use the committed corrected driver.

See \path{docs/performance/STEP4.md} for the declared gates and commands.
\path{benchmarks/gpu_batched_suite.py} reproduces the scaling/long-run job order;
\path{benchmarks/gpu_batched_quality.py} applies posterior gates with wall-time
ESS/s. Outputs must use new directories. GPU hardware tests require
\texttt{IMPALA\_TEST\_GPU=1}. Normal CPU tests skip optional GPU dependencies.

Source, tests, compact results, diagnostics and LaTeX source are committed
locally. Raw traces, logs and this PDF are in workspace
\path{benchmark/changes/step4/}; the versioned manifest records their hashes.
No changes have been published upstream. The next GPU performance hypothesis
would need to reduce host/transfer costs, not merely increase batch sizes.
No such redesign is claimed or implemented in this checkpoint.
\end{document}
